For an integer $n\ge 1$, define $k(n)$ to be the largest integer $k\ge 1$ such that there exists an integer $m\le n$ with the property that each of
\[
 m+1,\;m+2,\;\dots,\;m+k
\]
is divisible by at least one prime $p>k$.

Equivalently, for each $i\in\{1,\dots,k\}$, the integer $m+i$ has a prime factor exceeding $k$.

Edge case: $k(1)=1$ because taking $m=1$ gives $m+1=2$ divisible by a prime $2>1$, and $k\ge 2$ is impossible since $m+1=2$ is not divisible by a prime $>2$.
